\section{Implementation and Reproduction Details}
\label{sec:reproducibility}
\subsection{Evaluation Protocol and Runtime Settings}
The primary experiments were conducted on September 21--23, 2026. The mixed
workload contains 256 label-balanced BoolQ validation examples and five MMLU
test examples from each of 57 subjects, for 541 requests. Sampling uses seed
20260921. The RTX 4090 comparison uses a fixed 285-question MMLU subset.
Data preparation verifies source checksums. JET uses zero-shot semantic
candidate scoring; the separate manually reported Jev 1.13 evaluation uses the
complete MMLU test set (89.06\%) and reports 2.86 requests/s via API. Jev does not
expose configurable thinking or few-shot settings. API concurrency and detailed
timing boundaries are not available in the supplied measurement.

The mixed CPU/GPU experiment uses an i9-13900K, 64 GiB RAM, and a 16 GiB
Arc A770. Qwen3.6-35B-A3B Q4\_K\_M places the first 11 layers' routed experts
on CPU, using eight CPU threads, two sequence slots, 2,048 context tokens per
sequence, microbatch/output rows of 256, allocated weights
(\texttt{--no-mmap}), and disabled thinking. Each worker-lifetime variant has
one full run. A separate 32-request diagnostic probe is not an accuracy sample.

The RTX 4090 experiment uses full GPU placement, two sequence slots,
microbatch/output rows of 256, eight CPU threads, eight input requests per
chunk, allocated weights, and disabled thinking. Runs execute sequentially in
fresh processes; approximately ten seconds of loading are included. All three
variants execute 285 prefills, 60,680 prefill tokens, 9,283 continuation tokens,
and 1,472 decoder calls. Preparation overlaps scoring, so phase times are not
additive. Exploratory runs affected by concurrent compilation are excluded.

The prefix comparison uses the preserved \texttt{ac0477b} binary as baseline.
Both versions use eight input requests per chunk, nine sequence slots, 2,048
context tokens per sequence, token batch 2,048, microbatch 512, and output-row
limit 256. The baseline runs once per model; the new version runs twice, with
model order reversed in the second round and no explicit warm-up. New execution
uses 541 prefills and 128,583 prefill tokens; the baseline schedule implies
1,652 prefills and 378,526 tokens. Continuations comprise 9,795 tokens across
2,369 decoder calls. Optional thinking uses a 256-token quota; complete-process
time rises from 69.51 to 644.23 seconds and combined NLL from 2.0977 to 2.2883.

\subsection{Source and Model Versions}
The implementation description follows repository commit
\texttt{4ff0ab63018e} and llama.cpp submodule commit
\texttt{b29c606e28a0}. Benchmark binaries are separate artifacts;
these source identifiers do not imply that every reported experiment ran at
that revision. The experiment records provide individual dates, settings, commands, and
artifact locations. Internal records are indexed in \path{paper/evidence.md}.

The Qwen3.6 text model is the Q4\_K\_M conversion from
\texttt{ggml-org/Qwen3.6-35B-A3B-GGUF}, revision
\texttt{baec3ebee244827cda0f4557eafa8b28f7545fa6}.
Its file size is 20,419,565,568 bytes and its SHA-256 is
\par\noindent\begin{minipage}{\linewidth}
\begin{lstlisting}
671e47e0ec53c665d048b98c3ecbfd5236b5ca9c3e02ed19fc8f81f7b85140c7
\end{lstlisting}
\end{minipage}\par
The text benchmark does not load a vision projector. Model download scripts pin
and verify the remaining evaluation weights and the optional matching projector.

\subsection{Experimental Artifacts}
All paths below are relative to \texttt{tests/accuracy/generated/}. These
artifacts are excluded from Git; a source checkout alone does not contain them.
\begin{itemize}
  \item \path{vulkan-prefix-benchmark-20260922/}: old/new binary comparisons,
  response files, stage counters, and two new runs for each Qwen3.5 model.
  \item \path{qwen36-35b-a3b-q4-offload-20260922/}: original mixed CPU/Vulkan
  placement study, model provenance, and full-workload results.
  \item \path{qwen36-profile-20260922/full-pool541/}: persistent-worker
  full-workload run; its parent directory contains diagnostic and comparison data.
  \item \path{qwen36-preparation-20260923/}: RTX 4090 results.
  Run directories are \path{baseline-repeat/}, \path{final-cached-sync/}, and
  \path{final-pipeline/}. These paths are documented in the performance notes;
  they were not available in the local artifact set used to prepare this draft.
\end{itemize}
The prefix summaries and persistent-worker accuracy report were verified
against locally retained machine-readable artifacts. RTX 4090 measurements
are sourced from the versioned performance notes. The claim-to-source mapping
is provided in \path{paper/evidence.md}.

\subsection{Reproduction Commands}
The default sampled accuracy workflow is:
\par\noindent\begin{minipage}{\linewidth}
\begin{lstlisting}
./scripts/download-test-model.sh
JET_BACKEND=vulkan ./scripts/run-accuracy-eval.sh
\end{lstlisting}
\end{minipage}\par
It requires the Vulkan build dependencies documented in the repository and uses
the pinned Qwen3-0.6B fixture. After preparing the standard inputs and building
a release Vulkan binary, the recorded larger-model placement can be rerun with:
\par\noindent\begin{minipage}{\linewidth}
\begin{lstlisting}
./scripts/download-accuracy-models.sh --qwen36
python3 scripts/benchmark-model.py \
  --model-path models/Qwen3.6-35B-A3B-Q4_K_M.gguf \
  --model-id qwen/qwen3.6-35b-a3b-q4_k_m \
  --output-dir tests/accuracy/generated/report-rerun \
  --runs 1 \
  --extra-args --cpu-moe-layers 11 --threads 8 \
  --max-sequences 2 --micro-batch 256 \
  --max-output-rows 256 --no-mmap --thinking disabled
\end{lstlisting}
\end{minipage}\par
Use a new output directory for each configuration. Reproducing the
before/after results additionally requires the preserved binaries or their
corresponding source snapshots; running the current binary alone does not
recreate those baselines.

For new measurements, retain complete command lines, model and projector hashes,
dataset manifests, prompts and candidate order, seeds, hardware and driver
versions, build features, placement, batch settings, and execution mode.
Report warm-up and repetition policy explicitly. Keep complete scorer time
separate from its enclosed phases and account for overlap in preparation timing.
